\documentclass[11pt,a4paper]{article}

\usepackage[margin=1in]{geometry}
\usepackage{amsmath, amssymb, amsthm, mathtools}
\usepackage{bm}
\usepackage{hyperref}
\usepackage{enumitem}
\usepackage{booktabs}
\usepackage{xcolor}

\newtheorem{theorem}{Theorem}
\newtheorem{proposition}[theorem]{Proposition}
\newtheorem{lemma}[theorem]{Lemma}
\newtheorem{corollary}[theorem]{Corollary}
\theoremstyle{definition}
\newtheorem{definition}[theorem]{Definition}
\newtheorem{remark}[theorem]{Remark}

\newcommand{\rr}{\mathbf{r}}
\newcommand{\xx}{\mathbf{x}}
\newcommand{\hh}{\mathbf{h}}
\newcommand{\nh}{\hat{\mathbf{n}}}
\newcommand{\kh}{\hat{\mathbf{k}}}
\newcommand{\Sab}{S_{\mathrm{ab}}}
\newcommand{\Q}{\mathbf{Q}}
\newcommand{\G}{\mathbf{G}}
\newcommand{\Gt}{\tilde{\mathbf{G}}}
\newcommand{\F}{\mathbf{F}}
\newcommand{\Cp}{\mathbf{C}}
\newcommand{\psib}{\boldsymbol{\psi}}
\newcommand{\Sigmaset}{\Sigma}
\newcommand{\C}{\mathbb{C}}
\newcommand{\R}{\mathbb{R}}
\newcommand{\E}{\mathbb{E}}
\newcommand{\trace}{\operatorname{tr}}
\newcommand{\rank}{\operatorname{rank}}
\newcommand{\lammax}{\lambda_{\max}}
\newcommand{\diff}{\mathrm{d}}

\title{Direction 5 (math note): \\
       Identifiability of the exposure operator \texorpdfstring{$\Q$}{Q}
       from the UE channel \texorpdfstring{$\hh$}{h}}
\author{(working note in preparation for JSAC SI submission)}
\date{\today}

\begin{document}
\maketitle

\begin{abstract}
The exposure-constrained beamformer (ECBF) of the AEGIS monograph requires both
the UE channel vector $\hh\in\C^M$ and the body exposure operator $\Q\in\C^{M\times M}$.
Standard pilot sounding delivers $\hh$ but not $\Q$. This note formalises the
gap. We give an information--geometric argument that $\Q$ is generically
non-identifiable from $\hh$ alone, characterise the minimal additional
measurements that close the gap (formalising the three ``recovery routes''
of the working brief), and prove a first-order perturbation bound on $\Q$
under rigid body motion. The bound separates cleanly into a translation
term (sub-wavelength sensitivity) and a rotation term (radian sensitivity),
which leads directly to a two-timescale estimation architecture in which
$\hh$ is updated per coherence slot and $\Q$ per body-pose interval.
We close with an $\hh$-observable necessary condition for the
exposure-signal alignment $\rho$ to be large, which lets the BS decide
when ECBF is actually needed without knowing $\Q$ exactly.
\end{abstract}

\section{Setup and notation}
\label{sec:setup}

We adopt the conventions of Part~III of the AEGIS monograph
(\S\,Field channel matrix, \S\,Exposure operator).
A base station with $M$ antenna elements illuminates a scene containing one
or more bodies and one UE at $\rr_{\mathrm{UE}}$. Site-specific propagation
generates $N$ paths. Path $n$ has

\begin{itemize}[leftmargin=2em,itemsep=2pt]
\item arrival direction $\kh_n \in S^2$ at the body region of interest;
\item polarisation--amplitude vector
      $\psib_n \in \C^3$ with $\psib_n \perp \kh_n$
      (so $\psib_n$ has $4$ real degrees of freedom, encoded for instance
      in the $(\hat{\boldsymbol{\theta}}(\kh_n),\hat{\boldsymbol{\varphi}}(\kh_n))$
      basis);
\item originating BS element index $j(n)\in\{1,\dots,M\}$.
\end{itemize}

The UE has receive antenna pattern $\Cp_R(\kh)\in\C^2$. The
\emph{path data} is the multi-set
\[
  \mathcal{P} \;=\; \bigl\{(\psib_n,\kh_n,j(n))\bigr\}_{n=1}^N
  \;\in\;\bigl(T^*S^2\bigr)^N \times \{1,\dots,M\}^N,
\]
where $T^*S^2$ denotes the cotangent bundle of the unit sphere
(continuous part has $4+2=6$ real dof per path).

\paragraph{Channel projection.}
The UE channel vector
$\hh = (h_1,\dots,h_M)^T \in \C^M$ has entries
\begin{equation}
  h_j \;=\; \sum_{n:\,j(n)=j}
            a_n \, e^{-ik_0 \kh_n \cdot \rr_{\mathrm{UE}}},
  \qquad
  a_n \;=\; \Cp_R(\kh_n)^H \psib_n \;\in\;\C.
  \label{eq:h-def}
\end{equation}
We refer to $\pi_h:\mathcal{P}\to\C^M$ as the
\emph{channel projection} and to $a_n$ as the \emph{antenna-collapsed
amplitude} of path $n$.

\paragraph{Exposure operator.}
For a body surface $\Sigmaset$ with outward unit normal $\nh(\rr)$, define
the per-path Fresnel transmission operator
$\F_n(\rr): \C^3 \to \C^3$ as in
the monograph (TE/TM amplitude coefficients $t_{s,n}, t_{p,n}$ depending on
$\mu_n(\rr) = -\kh_n\cdot\nh(\rr)$, supported on the front-facing region
$\mu_n>0$). The body-surface field channel is
\begin{equation}
  \Gt_j(\rr) \;=\; \sum_{n:\,j(n)=j}
                   \F_n(\rr)\,\psib_n\,e^{-ik_0\kh_n\cdot\rr},
\end{equation}
and the exposure operator is the Gram integral
\begin{equation}
  \Q \;=\; \int_{\Sigmaset} \Gt(\rr)^H \Gt(\rr)\,\diff A
       \;\in\;\C^{M\times M}, \qquad
  P_{\mathrm{abs}} \;=\; \xx^H \Q \xx.
  \label{eq:Q-def}
\end{equation}
We refer to $\pi_Q:\mathcal{P}\to\mathcal{H}_M^+$ as the
\emph{operator construction}, where $\mathcal{H}_M^+$ is the cone of
Hermitian PSD $M\times M$ matrices.

\paragraph{ECBF.}
The closed-form ECBF precoder is
\begin{equation}
  \xx^\star \;=\; \sqrt{P}\,
                  \frac{(\lambda\Q + \nu \mathbf{I})^{-1}\hh^*}
                       {\|(\lambda\Q + \nu \mathbf{I})^{-1}\hh^*\|},
  \label{eq:ecbf}
\end{equation}
with multipliers $(\lambda,\nu)\ge 0$ from KKT. Both $\hh$ and $\Q$
are required at the BS.

\paragraph{Notation for perturbations.}
Body rigid motion is parameterised by $g=(R,\mathbf{t})\in SE(3)$ with
infinitesimal generators $(\delta\boldsymbol{\omega},\delta\mathbf{t})\in
\mathfrak{se}(3)$. We write $\Q(g)$ for the operator under pose $g$ and
linearise around the identity.

% =================================================
\section{The identifiability gap}
\label{sec:identifiability}

\subsection{A degree-of-freedom counting argument}

The path data $\mathcal{P}$ has $6N$ continuous real dof
(plus $N$ discrete element labels). The channel projection $\pi_h$
maps into $\C^M$, which has $2M$ real dof. The operator construction
$\pi_Q$ maps into $\mathcal{H}_M^+$, which has $M^2$ real dof.

\begin{proposition}[Per-path information loss]
\label{prop:per-path-loss}
For each path $n$, the channel projection retains
$2$ of the $4$ continuous polarisation dof of $\psib_n$ and $0$ of
the $2$ continuous arrival-angle dof of $\kh_n$ (only a phasor
$e^{-ik_0\kh_n\cdot\rr_{\mathrm{UE}}}$ is retained, which couples to
$\kh_n$ only at the single point $\rr_{\mathrm{UE}}$). The operator
construction retains all $6$ continuous dof of every path.
\end{proposition}

\begin{proof}
The per-path channel contribution
$a_n = \Cp_R(\kh_n)^H\psib_n$ is a rank-$1$ linear functional on
$\psib_n\in\C^2 \cong \C^2$ (using the
$(\hat{\boldsymbol{\theta}},\hat{\boldsymbol{\varphi}})$ basis at $\kh_n$),
giving $2$ real dof out of $4$. The arrival direction enters $h_j$ only
through the scalar phase $e^{-ik_0\kh_n\cdot\rr_{\mathrm{UE}}}$;
two paths with distinct $\kh_n$ that share this phase are
indistinguishable in $h_j$ at the single point $\rr_{\mathrm{UE}}$.

For $\pi_Q$: the per-path Fresnel operator
$\F_n(\rr)\psib_n e^{-ik_0\kh_n\cdot\rr}$ depends on $\psib_n$ as a
rank-$2$ map ($\F_n$ injective on its range)
and on $\kh_n$ through both $\mu_n(\rr)$ and the phase ramp
$e^{-ik_0\kh_n\cdot\rr}$ over the surface.
\end{proof}

\begin{corollary}[Structural unidentifiability]
\label{cor:structural-unid}
Whenever $N\ge M$, the map $\pi_Q$ does not factor through $\pi_h$: there
exist $\mathcal{P},\mathcal{P}'$ with
$\pi_h(\mathcal{P})=\pi_h(\mathcal{P}')$ but
$\pi_Q(\mathcal{P})\ne\pi_Q(\mathcal{P}')$.
\end{corollary}

\begin{proof}
Fix any $\mathcal{P}$ with $h_j$ as in \eqref{eq:h-def}.
For any path $n$, replace $\psib_n$ by
\[
  \psib'_n \;=\; \psib_n + c\,\boldsymbol{\eta}_n,
\]
where $\boldsymbol{\eta}_n\in\C^3$ is non-zero, satisfies
$\boldsymbol{\eta}_n\perp\kh_n$, and lies in the kernel of
$\Cp_R(\kh_n)^H$ (this kernel is one-complex-dimensional, hence non-trivial,
because $\Cp_R^H$ is a linear map from $\C^2\cong\C^2$ to $\C$).
The channel value $a_n=\Cp_R(\kh_n)^H\psib_n$ is unchanged, so $\hh$ is
unchanged. But $\F_n\boldsymbol{\eta}_n\ne 0$ generically, so
$\Gt$, $\tilde{\Gt}$, and hence $\Q$ change.
\end{proof}

The kernel direction $\boldsymbol{\eta}_n$ is the
\emph{cross-polarisation} component invisible to the UE antenna pattern but
visible to the body surface (which absorbs both polarisations weighted by
TE/TM Fresnel coefficients). Corollary~\ref{cor:structural-unid} therefore
identifies the structural reason ECBF cannot be closed by uplink sounding
alone: the UE antenna is a \emph{polarisation filter}, and $\Q$ depends on
the cross-polarisation component that this filter discards.

\subsection{What a sufficient statistic for $\Q$ looks like}

\begin{proposition}[Sufficient statistic]
\label{prop:sufficient}
Given the body geometry $(\Sigmaset,\nh,\sigma,\rho_{\mathrm{tissue}},\tilde n)$,
$\Q$ is a deterministic functional of the body-surface field channel
$\Gt:\Sigmaset\to\C^{3\times M}$. Equivalently, $\Q$ is a deterministic
functional of the path data $\{(\psib_n,\kh_n,j(n))\}_{n=1}^N$ together
with the body geometry.
\end{proposition}

This is immediate from \eqref{eq:Q-def} but worth stating: $\Q$ has no
``hidden'' dependence on traffic, on noise, or on temporal averaging.
Whatever recovers $\Gt$ recovers $\Q$, and whatever recovers
$\{(\psib_n,\kh_n)\}$ together with body geometry recovers $\Gt$.
The operational question is which of these is cheapest to acquire.

% =================================================
\section{The three recovery routes, formally}
\label{sec:routes}

\subsection{Route A: ray tracing on a digital twin}

Assume the BS holds a geometric model $\mathcal{S}$ of the scene
(walls, scatterers, body pose at known location). A differentiable ray
tracer $\mathcal{R}:\mathcal{S}\to\mathcal{P}$ delivers the path data,
and $\Q = \pi_Q(\mathcal{R}(\mathcal{S}),\Sigmaset)$ follows from
\eqref{eq:Q-def}. The DiffeRT integration in AEGIS realises this map.

\paragraph{Sufficiency.} Trivially sufficient when $\mathcal{S}$ matches
ground truth: $\pi_Q$ is by definition computable from $\mathcal{P}$.

\paragraph{Failure modes.} The error
$\Q-\Q^{true}$ is driven by mismatch between
$\mathcal{S}$ and the true scene. Section~\ref{sec:perturbation} bounds
this for the dominant component (body pose mismatch).

\subsection{Route B: pilot sounding sufficient to recover \texorpdfstring{$\psib_n$}{psi\_n}}

To upgrade pilots from delivering $\hh$ to delivering $\{\psib_n\}$ we need
to invert the per-path projection
$a_n=\Cp_R(\kh_n)^H\psib_n$. Because $\Cp_R(\kh_n)^H:\C^2\to\C^1$ is
rank-$1$, a single scalar $a_n$ cannot determine $\psib_n$.

\begin{proposition}[Minimal pilot diversity for $\psib_n$]
\label{prop:pilot-diversity}
Suppose the UE has a dual-polarised receive front end with two
linearly independent patterns
$\Cp_R^{(1)}(\kh),\Cp_R^{(2)}(\kh)\in\C^2$, and that the BS can
estimate the per-path scalar from each port. Suppose further that the
arrival direction $\kh_n$ of each path is known (e.g.\ via UE-side AoA
estimation from a small antenna array). Then $\psib_n$ is identifiable.
\end{proposition}

\begin{proof}
With $\kh_n$ known, $\psib_n$ lives in the $2$-complex-dimensional plane
$\kh_n^{\perp}\subset\C^3$. The two scalar measurements
\[
  a_n^{(p)} \;=\; \Cp_R^{(p)}(\kh_n)^H\psib_n,\qquad p=1,2,
\]
form a $2\times 2$ linear system on $\psib_n$ in the
$(\hat{\boldsymbol{\theta}},\hat{\boldsymbol{\varphi}})$ basis at $\kh_n$.
Linear independence of $\Cp_R^{(1)},\Cp_R^{(2)}$ guarantees
invertibility.
\end{proof}

The practical price is dual-polarised UE Rx hardware (already standard for FR2),
UE-side AoA estimation (already done for spatial multiplexing), and a
small protocol extension that exposes per-path polarisation amplitudes
to the BS rather than the collapsed channel. Note this still does not
deliver $\Q$ directly; it delivers the path-level data needed to
\emph{construct} $\Q$ given a body model. So Route~B reduces to
Route~A with a different source of $\mathcal{P}$ (sounded paths instead of
ray-traced paths) and the same dependence on a body-pose model.

\subsection{Route C: phantom-calibrated \texorpdfstring{$\Q$}{Q} envelope}

Hochwald et al.\ measure their SAR matrix $\mathbf{S}_V$ on a calibrated
phantom in a fixed canonical pose. The same procedure applied to AEGIS
gives a phantom-calibrated $\Q^{\mathrm{phantom}}$ that bounds absorbed
power from above:
\[
  P_{\mathrm{abs}}^{(u)} \;\le\; \lammax(\Q^{\mathrm{phantom}})\,P
  \quad\text{for every precoder with}\quad \|\xx\|^2\le P.
\]
This is a one-shot \emph{static} bound; it is independent of $\hh$ and so
yields a worst-case ECBF: replace $\Q$ in \eqref{eq:ecbf} by
$\Q^{\mathrm{phantom}}$. Section~\ref{sec:rho-witness} shows that this
bound is loose precisely in the regime where the $\hh$-witness is small,
which suggests pairing Routes~A/C: use Route~A's posterior $\Q$
when the $\hh$-witness flags ECBF as needed, and fall back to the static
Route~C bound otherwise.

% =================================================
\section{Perturbation theorem for \texorpdfstring{$\Q$}{Q} under rigid body motion}
\label{sec:perturbation}

We now bound $\Q(g) - \Q(0)$ for small $g=(R,\mathbf{t})\in SE(3)$.
Throughout this section the path data is treated as fixed (the scene
outside the body is unchanged) and $\Q$ varies because $\Sigmaset$,
$\nh$, and $\mu_n$ depend on body pose.

\subsection{Translation: sub-wavelength sensitivity}

\begin{proposition}[Translation perturbation]
\label{prop:translation}
Let $\Q(\mathbf{t})$ be the exposure operator after rigid translation of
the body by $\mathbf{t}\in\R^3$. The $(a,b)$ entry decomposes
exactly as
\begin{equation}
  [\Q(\mathbf{t})]_{ab} \;=\;
  \sum_{\substack{n:\,j(n)=a\\ n':\,j(n')=b}}
  e^{-ik_0(\kh_n - \kh_{n'})\cdot\mathbf{t}}\,
  Q_{ab}^{(nn')},
  \label{eq:Q-trans}
\end{equation}
where $Q_{ab}^{(nn')}$ is the path-pair contribution to $[\Q(0)]_{ab}$.
For $\|\mathbf{t}\|\ll \lambda_0/(2\pi)$ this gives the first-order
expansion
\begin{equation}
  \Q(\mathbf{t}) \;=\; \Q(0) - i k_0\,\mathbf{t}\cdot\Delta\Q \;+\; O\bigl((k_0\|\mathbf{t}\|)^2\bigr),
\end{equation}
with $\|\Delta\Q\|_F \le 2\|\Q(0)\|_F$, hence
\begin{equation}
  \frac{\|\Q(\mathbf{t})-\Q(0)\|_F}{\|\Q(0)\|_F}
  \;\le\; 2\,k_0\|\mathbf{t}\| \;+\; O\bigl((k_0\|\mathbf{t}\|)^2\bigr).
  \label{eq:Q-trans-bound}
\end{equation}
\end{proposition}

\begin{proof}[Sketch]
Translating the body by $\mathbf{t}$ shifts each surface point
$\rr\mapsto \rr+\mathbf{t}$ but leaves $\nh(\rr)$, $\mu_n$, and $\F_n$
invariant (Fresnel coefficients are rotation-only sensitive). Each
per-path field then acquires a phasor $e^{-ik_0\kh_n\cdot\mathbf{t}}$,
which factors out of the area integral and survives in the $\Gt^H\Gt$
quadratic as $e^{-ik_0(\kh_n-\kh_{n'})\cdot\mathbf{t}}$. Expanding the
exponential to first order and bounding by the sum of moduli gives the
Frobenius bound.
\end{proof}

\paragraph{Operational consequence.} At $28\,$GHz, $\lambda_0\!\approx\!1.07\,$cm,
so $k_0=2\pi/\lambda_0 \approx 587\,\mathrm{m}^{-1}$. A body translation
of $\|\mathbf{t}\|=1\,$mm gives $k_0\|\mathbf{t}\|\approx 0.59\,$rad,
which by \eqref{eq:Q-trans-bound} can move $\Q$ by more than its own
norm. \emph{Translation invalidates $\Q$ on a sub-wavelength scale.}

The bound is conservative because \emph{eigenstructure} of $\Q$ is more
robust than its Frobenius norm. We turn to that next.

\begin{corollary}[Spectral robustness]
\label{cor:spectral}
By Weyl's inequality,
$|\lammax(\Q(\mathbf{t})) - \lammax(\Q(0))|
 \le \|\Q(\mathbf{t})-\Q(0)\|_2 \le \|\Q(\mathbf{t})-\Q(0)\|_F$,
so the Frobenius bound transfers, but is only tight when the perturbation
is rank-one along the leading eigenvector. In practice $\delta\Q$ from
translation is approximately a path-pair phase rotation, which mixes
many eigenmodes and therefore typically affects $\lammax$ less than the
Frobenius bound suggests.
\end{corollary}

\subsection{Rotation: radian sensitivity}

\begin{proposition}[Rotation perturbation]
\label{prop:rotation}
Let $\Q(\theta\hat{\boldsymbol{\omega}})$ denote the exposure operator
after rigid rotation of the body by angle $\theta$ about axis
$\hat{\boldsymbol{\omega}}\in S^2$. Assume the visible-path set is
unchanged (no self-shadowing transitions). Then
\begin{equation}
  \frac{\|\Q(\theta\hat{\boldsymbol{\omega}})-\Q(0)\|_F}{\|\Q(0)\|_F}
  \;\le\; C_{\mathrm{rot}}\,\theta \;+\; O(\theta^2),
\end{equation}
where $C_{\mathrm{rot}}$ depends on
(i) the maximum slope of the Fresnel amplitude coefficients
$|\partial_\mu t_{s,n}|$, $|\partial_\mu t_{p,n}|$ over visible paths
and (ii) the diameter of the body, but not on $k_0$.
\end{proposition}

\begin{proof}[Sketch]
Rigid rotation perturbs the per-point normal as
$\nh(\rr;\theta)=\nh(\rr;0)+\theta(\hat{\boldsymbol{\omega}}\times\nh(\rr))+O(\theta^2)$,
so $\delta\mu_n(\rr)=-\theta\,(\nh\times\kh_n)\cdot\hat{\boldsymbol{\omega}} = O(\theta)$.
The Fresnel amplitudes are smooth in $\mu$ on $\mu>0$ (the only loss of
smoothness is the $\mu_n=0$ horizon, which by assumption no path crosses).
The TE/TM basis vectors rotate with $\nh$ at the same rate.
The phase ramp $e^{-ik_0\kh_n\cdot\rr}$ acquires a per-point variation
$e^{-ik_0\kh_n\cdot\delta\rr(\theta)}$ where $\delta\rr=O(\theta\,\|\rr\|)$;
this contributes a body-diameter-scaled rotation term. Combining and
re-integrating over $\Sigmaset$ gives the stated bound.
\end{proof}

\paragraph{Operational consequence.} Unlike translation, rotation does
\emph{not} carry a $k_0$ multiplier in front of the smooth-Fresnel
component. Body rotation by $\theta\!=\!0.1\,$rad ($\sim 6^\circ$)
typically perturbs $\Q$ by $\sim\!10\%$ in Frobenius norm. The
phase-ramp term, however, does carry a $k_0\!\cdot\!\mathrm{diam}(\Sigmaset)$
factor; for a body of diameter $50\,$cm at $28\,$GHz this factor is
$\sim\!300$, so even small rotations can scramble inter-path interference
patterns far away from the rotation axis. The bound is again conservative
on $\lammax$.

\subsection{Self-shadowing transitions}

The smooth bounds above assume the visible-path set is unchanged.
A path with $\mu_n(\rr;0)=0$ at some $\rr$ enters or leaves the
visible set under arbitrarily small pose perturbation, with discontinuous
contribution to $\Q$. These transitions are \emph{measure zero} in pose
space (a finite union of codimension-1 strata) but their
\emph{cumulative} effect over a finite pose excursion is bounded by the
total grazing-incidence flux:
\[
  \|\delta\Q^{\mathrm{shadow}}\|_F \;\le\;
  \int_{\partial\Sigmaset_{\mathrm{vis}}}
  \left\|\Gt^{(n)}(\rr)\right\|^2 \diff\ell,
\]
i.e.\ a line integral along the silhouette curve weighted by the
grazing-path field. This is small whenever no dominant path is at
near-grazing incidence on the body, which is the generic case for
ECBF-relevant scenarios (the dominant exposure modes correspond to
near-normal incidence).

% =================================================
\section{Two-timescale architecture}
\label{sec:timescales}

The perturbation analysis motivates a decoupled estimation architecture.
Let $T_c$ be the channel coherence time (typical $1{-}10\,$ms at FR2 for
mobile users) and $T_p$ the body-pose coherence time (typical
$0.1{-}1\,$s for a quasi-stationary user; longer for static seated/standing
postures).

\begin{definition}[Two-timescale ECBF]
\label{def:two-timescale}
Pilot sounding refreshes $\hh(t)$ at every slot of length $\le T_c$.
A separate body-pose channel refreshes $\Q(t_p)$ at intervals of
length $\le T_p$. The per-slot precoder is
\[
  \xx^\star(t) \;=\; \sqrt{P}\,
                  \frac{(\lambda(t)\Q(t_p^-) + \nu(t)\mathbf{I})^{-1}\hh(t)^*}
                       {\|\cdot\|},
\]
with multipliers re-solved per slot from current $(\hh(t),\Q(t_p^-))$,
where $t_p^-$ is the most recent body-pose update.
\end{definition}

\begin{proposition}[Time-averaged constraint under staleness]
\label{prop:time-avg}
Let $P_{\mathrm{abs}}^{(true)}(t)$ be the actual absorbed power
under $\xx^\star(t)$ and $\Q^{(true)}(t)$, and let
$P_{\mathrm{abs}}^{(des)}(t)$ be the value computed at design time
using $\Q(t_p^-)$. Then over a pose interval $[t_p^-,t_p^-+T_p)$,
\[
  \left|\frac{1}{T_p}\int_{t_p^-}^{t_p^-+T_p}
        \!\!\bigl(P_{\mathrm{abs}}^{(true)}(t) - P_{\mathrm{abs}}^{(des)}(t)\bigr)\diff t\right|
  \;\le\; P\,\overline{\|\delta\Q\|_2}_{[t_p^-,t_p^-+T_p)},
\]
where $\overline{\|\delta\Q\|_2}$ is the time-averaged spectral
distance between $\Q^{(true)}$ and $\Q(t_p^-)$ over the interval.
By Propositions~\ref{prop:translation}--\ref{prop:rotation}, this is
controlled by the time-averaged body translation/rotation since the last
pose update.
\end{proposition}

\begin{proof}
Given any unit-norm $\xx$, $|\xx^H(\Q^{(true)}-\Q(t_p^-))\xx|\le
\|\delta\Q\|_2$. With $\|\xx^\star(t)\|^2=P$, multiply by $P$ and
average in $t$.
\end{proof}

This gives a clean separation of concerns. The compliance margin one
must reserve for staleness scales linearly with body-motion budget,
not with channel dynamics.

% =================================================
\section{An \texorpdfstring{$\hh$}{h}-observable witness for ECBF necessity}
\label{sec:rho-witness}

The exposure-signal alignment $\rho$ from the monograph quantifies how much
of the worst-case absorption MRT actually achieves:
\begin{equation}
  \rho \;=\; \frac{\hh^T \Q\,\hh^*}{\|\hh\|^2\,\lammax(\Q)} \;\in\; [0,1].
  \label{eq:rho-def-restate}
\end{equation}
The BS would like to decide ``ECBF is unnecessary; MRT will satisfy the
exposure budget'' from $\hh$ alone. This is impossible without
\emph{some} information about $\Q$, but it is possible to bound $\rho$
from above using only $\hh$ together with a coarse \emph{angular
support} $K\subset S^2$ of body-incident arrival directions.

\begin{proposition}[\texorpdfstring{$\hh$}{h}-witness for low \texorpdfstring{$\rho$}{rho}]
\label{prop:rho-witness}
Suppose the BS has access to (i) the channel $\hh$, (ii) the per-path
arrival directions $\{\kh_n\}_{n=1}^N$ at the BS array (recoverable
from spatial covariance via MUSIC/ESPRIT in the multi-snapshot regime),
and (iii) a coarse angular set $K\subset S^2$ that contains all
directions illuminating the body (e.g.\ a cone subtended by the body's
silhouette as seen from the BS). Decompose
\[
  \hh \;=\; \hh_K + \hh_{K^c},
  \qquad
  \hh_K \;=\; \sum_{n:\,\kh_n\in K} a_n e^{-ik_0\kh_n\cdot\rr_{\mathrm{UE}}}\mathbf{e}_{j(n)}.
\]
Then
\begin{equation}
  \rho \;\le\; \frac{\|\hh_K\|^2}{\|\hh\|^2}.
  \label{eq:rho-witness}
\end{equation}
\end{proposition}

\begin{proof}
The exposure operator $\Q$ has support entirely in the body-illuminating
subspace, in the sense that $\Q\,\hh_{K^c}^* = \mathbf{0}$ (every path
contributing to $\Gt$ has $\kh_n\in K$). Hence
$\hh^T\Q\hh^* = \hh_K^T\Q\hh_K^*$. By Rayleigh--Ritz on the PSD
matrix $\Q$, $\hh_K^T\Q\hh_K^*\le\lammax(\Q)\|\hh_K\|^2$.
Dividing by $\|\hh\|^2\,\lammax(\Q)$ yields the bound.
\end{proof}

\begin{remark}[Operational reading]
\label{rem:operational}
The right-hand side of \eqref{eq:rho-witness} is the
\emph{body-incident energy fraction} of the channel. It is computable
from $\hh$ and BS-side AoA estimates plus a coarse body location. If
this fraction is below some threshold $\epsilon$, then $\rho\le\epsilon$
unconditionally on the body's exact shape and pose, so MRT loses at most
a fraction $\epsilon$ of $\lammax(\Q)\,P$ to the body. Combined with
the static Route~C envelope $\lammax(\Q)\le\Lambda$, this gives an
$\hh$-only certificate that MRT is exposure-safe whenever
$\epsilon\Lambda\le P_{\mathrm{abs}}^{\max}/P$.
\end{remark}

\begin{remark}[What the witness does \emph{not} say]
The converse fails. Large $\|\hh_K\|^2/\|\hh\|^2$ does not imply $\rho$
is large; the body might still be illuminated mostly along low-Fresnel-
coupling directions, or the dominant body modes might be misaligned with
the channel within the body-incident subspace. The witness is a
\emph{necessary} condition for ECBF to matter, not sufficient. Sufficiency
requires Q.
\end{remark}

% =================================================
\section{What this lets us write in the JSAC SI submission}
\label{sec:writeup}

The above isolates four claims that can be made cleanly without any
additional empirical work, and three follow-up claims that require
small-scale numerical experiments in AEGIS:

\paragraph{Closed-form claims (this note alone).}
\begin{enumerate}[leftmargin=2em,itemsep=2pt]
\item Structural unidentifiability of $\Q$ from $\hh$
      (Cor.~\ref{cor:structural-unid}). This is the analogue, in the
      coherent-MIMO setting, of the polarisation-loss observation that
      already appears in the monograph (\S Signal branch); we sharpen it
      to a non-identifiability statement.
\item Minimal pilot diversity for $\psib_n$ recovery
      (Prop.~\ref{prop:pilot-diversity}). Identifies the minimum
      hardware/protocol upgrade that closes the gap.
\item Translation--rotation perturbation bounds
      (Props.~\ref{prop:translation}--\ref{prop:rotation}). Quantifies the
      tolerable body motion before $\Q$ must be re-estimated.
\item $\hh$-witness for ECBF necessity (Prop.~\ref{prop:rho-witness}).
      Lets the BS gate ECBF on a $\hh$-only observable.
\end{enumerate}

\paragraph{Empirical follow-ups (require AEGIS runs).}
\begin{enumerate}[leftmargin=2em,itemsep=2pt,start=5]
\item \textbf{Tightness of perturbation bounds.}
      For a representative scenario (one phantom, one BS, single-path or
      few-path environment, $28\,$GHz), compute
      $\|\Q(g)-\Q(0)\|_F/\|\Q(0)\|_F$ and
      $|\lammax(\Q(g))-\lammax(\Q(0))|/\lammax(\Q(0))$
      as a function of body translation in $\{0,\,0.1,\,0.5,\,1,\,5,\,10\}\,$mm
      and rotation in $\{0,\,1,\,5,\,10,\,30\}^\circ$. Compare to
      Prop.~\ref{prop:translation}--\ref{prop:rotation} bounds.
      Hypothesised result: Frobenius bounds are tight; spectral bounds
      have an order-of-magnitude slack.
\item \textbf{Witness sharpness.}
      For a sweep over UE positions and body locations, compare the
      $\hh$-witness $\|\hh_K\|^2/\|\hh\|^2$ to the true $\rho$. Identify
      the regime where the witness is informative (gap less than 1\,dB).
\item \textbf{Two-timescale ECBF in simulation.}
      Implement the protocol of Def.~\ref{def:two-timescale} on a
      synthetic body-motion trace and compare end-to-end exposure
      compliance to (a) per-slot oracle ECBF and (b) static
      Hochwald-style worst-case ECBF. Quantify the spectral-efficiency
      gain at fixed compliance probability.
\end{enumerate}

A clean way to structure the JSAC SI manuscript is to lead with the
identifiability and witness results (theory), then present the
two-timescale ECBF as the algorithmic contribution, and use AEGIS as
the digital-twin substrate that makes Route~A computationally feasible
at the per-slot cadence. The resulting story sits squarely in the
``twin-in-the-loop'' / ``closed-loop optimisation'' theme of the call.

% =================================================
\appendix

\section{Symbol table}
\label{app:symbols}

\begin{center}
\small
\begin{tabular}{lp{8cm}l}
\toprule
Symbol & Meaning & Space \\
\midrule
$M$ & number of BS antenna elements & scalar \\
$N$ & number of propagation paths & scalar \\
$\rr_{\mathrm{UE}}$ & UE position & $\R^3$ \\
$\Sigmaset$ & body surface & 2-manifold in $\R^3$ \\
$\nh(\rr)$ & outward unit normal at $\rr\in\Sigmaset$ & $S^2$ \\
$\kh_n$ & arrival direction of path $n$ at the body region & $S^2$ \\
$\psib_n$ & polarisation-amplitude of path $n$, $\psib_n\perp\kh_n$ & $\C^3$ \\
$j(n)$ & BS element originating path $n$ & $\{1,\dots,M\}$ \\
$a_n$ & UE-projected scalar amplitude $\Cp_R^H\psib_n$ & $\C$ \\
$\Cp_R(\kh)$ & UE receive antenna pattern & $\C^2$ \\
$\hh$ & UE channel vector & $\C^M$ \\
$\F_n(\rr)$ & per-path Fresnel transmission operator & $\C^{3\times 3}$ \\
$\Gt(\rr)$ & body-surface field channel matrix & $\C^{3\times M}$ \\
$\Q$ & exposure operator $\int_\Sigmaset \Gt^H\Gt\,\diff A$ & $\mathcal{H}_M^+$ \\
$\rho$ & exposure-signal alignment $\hh^T\Q\hh^*/(\|\hh\|^2\lammax(\Q))$ & $[0,1]$ \\
$g=(R,\mathbf{t})$ & body rigid pose & $SE(3)$ \\
$T_c, T_p$ & channel and pose coherence times & seconds \\
$K\subset S^2$ & angular support of body-incident directions at the BS & cone \\
$\hh_K, \hh_{K^c}$ & body-incident and non-body channel components & $\C^M$ \\
\bottomrule
\end{tabular}
\end{center}

\section{Cross-references to the monograph}
\label{app:xrefs}

\begin{center}
\small
\begin{tabular}{ll}
\toprule
This note & Monograph (\texttt{theory/monograph\_v2.tex}) \\
\midrule
\eqref{eq:h-def} & \texttt{eq:h-def}, \S\,Signal branch \\
\eqref{eq:Q-def} & \texttt{eq:Q-def}, Def.\,Q \\
\eqref{eq:ecbf} & \texttt{eq:optimal-x}, \S\,Optimal precoder \\
\eqref{eq:rho-def-restate} & \texttt{eq:rho-def}, Def.\,$\rho$ \\
Prop.~\ref{prop:per-path-loss} & paragraph after \texttt{eq:h-def} (informal) \\
Cor.~\ref{cor:structural-unid} & remark in \S\,Per-user exposure tracking \\
Routes A/B/C & remark in \S\,Per-user exposure tracking \\
\bottomrule
\end{tabular}
\end{center}

\end{document}
